\documentclass[../DoAn.tex]{subfiles}
\begin{document}

Chương này là chương trọng tâm của đồ án. Trên cơ sở các đặc tả ở Chương
\ref{section:2.3} và bộ công nghệ đã trình bày ở Chương \ref{section:3.1},
chương trình bày toàn bộ quá trình thiết kế, hiện thực hóa và đánh giá hệ
thống WAF. Nội dung gồm bảy phần: kiến trúc tổng quan và luồng xử lý request
xuyên suốt các lớp middleware (\ref{section:4.1}); thiết kế chi tiết rule
engine với cấu trúc dữ liệu, transform chain và cơ chế scoring
(\ref{section:4.2}); tích hợp ML layer và pipeline huấn luyện kiểm soát chất
lượng (\ref{section:4.3}); thiết kế cơ sở dữ liệu (\ref{section:4.4}); xây
dựng dashboard và các trang quản trị (\ref{section:4.5}); kết quả kiểm thử
toàn diện trên các tập payload tấn công chuẩn, kiểm thử tải và kiểm thử bảo
mật nội bộ (\ref{section:4.6}); và hướng dẫn triển khai thực tế qua Docker
Compose (\ref{section:4.7}).

% ====================================================================
\section{Kiến trúc tổng quan hệ thống}
\label{section:4.1}

\subsection{Sơ đồ kiến trúc tổng thể}
\label{subsection:4.1.1}

Hệ thống được tổ chức theo mô hình \textit{reverse proxy trong suốt}: client
gửi request đến WAF tại cổng công khai, WAF thực hiện toàn bộ kiểm tra, và
chuyển tiếp request đến backend nếu được phép. Backend không cần biết đến
sự tồn tại của WAF, không cần thay đổi cấu hình hay mã nguồn.

\begin{figure}[h!]
\centering
\includegraphics[width=0.55\textwidth]{Hinhve/architecture.png}
\caption{Sơ đồ kiến trúc tổng thể của hệ thống WAF}
\label{fig:architecture-overview}
\end{figure}

Hệ thống có ba hot path chính: \textbf{(i)} luồng xử lý request thông
thường (client \(\to\) WAF \(\to\) backend), \textbf{(ii)} luồng quản trị
(admin \(\to\) Dashboard \(\to\) WAF API \(\to\) PostgreSQL), và
\textbf{(iii)} luồng ML xác minh (rule engine \(\to\) ML Service \(\to\)
quay lại rule engine), được kích hoạt có điều kiện.

\subsection{Luồng xử lý request}
\label{subsection:4.1.2}

Một request HTTP đi qua chín bước xử lý tuần tự trong WAF core. Mỗi bước
có thể từ chối request và kết thúc sớm; nếu vượt qua tất cả, request được
forward đến backend.

\begin{enumerate}[leftmargin=*]
    \item \textbf{Nhận và parse request.} TCP connection được terminate, TLS
    được giải mã (nếu HTTPS), parser bóc tách method, path, query parameters,
    headers, và body với giới hạn kích thước \texttt{max\_body\_size}
    (mặc định 10 MB).

    \item \textbf{Normalize.} Áp dụng các biến đổi: URL decode một lớp,
    chuyển về lowercase, compress whitespace. Đầu ra là một \texttt{ParsedRequest}
    chuẩn hóa sẵn sàng cho các bước sau.

    \item \textbf{IP check.} Truy vấn cache trong bộ nhớ (đồng bộ định kỳ
    với PostgreSQL): nếu IP nguồn nằm trong blacklist, từ chối ngay với mã
    403; nếu nằm trong whitelist, bỏ qua các bước rule engine.

    \item \textbf{Rate Limit.} Token Bucket per-IP (mục \ref{section:3.7}):
    lấy đi một token; nếu bucket trống, từ chối với mã 429.

    \item \textbf{Authentication.} Với các endpoint quản trị của WAF
    (\texttt{/waf-api/*}, trang \texttt{/dashboard}), middleware đọc JWT từ
    cookie hoặc \texttt{Authorization} header; các trang công khai
    \texttt{/login.html}, \texttt{/register.html} không yêu cầu JWT. Toàn bộ
    bề mặt quản trị còn được bao bởi lớp \textit{admin access control}: IP
    ngoài \texttt{admin.allowed\_cidrs} (mặc định loopback) nhận 404 phẳng,
    không lộ sự tồn tại của giao diện quản trị.

    \item \textbf{Authorization.} Đối với các endpoint admin-only, kiểm
    tra claim \texttt{role} của JWT. Endpoint user-level chỉ yêu cầu JWT
    hợp lệ bất kể role.

    \item \textbf{Rule Engine.} Duyệt tuần tự 78 rule. Mỗi rule khớp đóng
    góp điểm vào tổng \texttt{anomaly\_score}. Bộ chọn rule (selector) loại
    bỏ sớm các rule không thể áp dụng cho request hiện tại (ví dụ rule body
    bỏ qua khi method GET).

    \item \textbf{ML Confirm (có điều kiện).} Nếu điểm anomaly nằm trong
    vùng xám \([3.0, 5.0)\), gọi ML service với canonical text, nhận kết
    quả, điều chỉnh điểm. Timeout cứng 800 mili giây, fail-open.

    \item \textbf{Decision.} So sánh điểm cuối với các ngưỡng (mặc định
    log \(\ge 3.0\), challenge \(\ge 4.0\), block \(\ge 5.0\)). Quyết định:
    BLOCK trả 403, CHALLENGE trả 429 hoặc 401 (tùy cấu hình), LOG/ALLOW
    forward đến backend.
\end{enumerate}

Mỗi bước ghi metric Prometheus (\texttt{waf\_requests\_total},
\texttt{waf\_decisions\_total}, \texttt{waf\_rule\_hits\_total},
\texttt{waf\_ml\_calls\_total}, \texttt{waf\_latency\_seconds}) cho phép
quan sát thông qua scrape từ Prometheus server.

Biểu đồ tuần tự (sequence diagram) ở Hình \ref{fig:seq-request-lifecycle}
biểu diễn các tương tác giữa các thành phần kể trên trong một vòng đời
request từ lúc client gửi đến lúc nhận phản hồi.

\begin{figure}[h!]
\centering
\includegraphics[width=0.95\textwidth,keepaspectratio]{Hinhve/seq_request_lifecycle.png}
\caption{Biểu đồ tuần tự -- Vòng đời một HTTP request đi qua WAF}
\label{fig:seq-request-lifecycle}
\end{figure}

\subsection{Các thành phần chính và tương tác}
\label{subsection:4.1.3}

Mã nguồn được tổ chức thành 14 package độc lập trong thư mục
\texttt{internal/}, mỗi package có một trách nhiệm hẹp và một interface
public rõ ràng. Bảng \ref{table:packages} liệt kê các package chính.

\begin{table}[h!]
\centering
\caption{Các package Go chính và trách nhiệm}
\label{table:packages}
\small
\begin{tabular}{|l|p{8.5cm}|}
\hline
\textbf{Package} & \textbf{Trách nhiệm} \\
\hline
\texttt{parser} & Bóc tách HTTP request thành \texttt{ParsedRequest}, giới hạn body size \\
\hline
\texttt{normalizer} & Áp dụng transform chain trên các trường của request \\
\hline
\texttt{engine} & Rule engine, transform, pattern matching, anomaly scoring, ML hook \\
\hline
\texttt{ml} & HTTP client gọi ML Inference Service, tuân thủ interface \texttt{MLPredictor} \\
\hline
\texttt{decision} & Áp dụng ngưỡng và quyết định ALLOW/LOG/CHALLENGE/BLOCK \\
\hline
\texttt{ratelimit} & Token Bucket per-IP với map đồng thời và GC định kỳ \\
\hline
\texttt{behavior} & Phát hiện hành vi bất thường (brute force, scanner, velocity) \\
\hline
\texttt{auth} & JWT generator, validator, RBAC middleware \\
\hline
\texttt{middleware} & Pipeline middleware: IP-check, rate limit, auth, logging \\
\hline
\texttt{api} & HTTP handler cho \texttt{/waf-api/*}: auth, users, rules, IPs, settings \\
\hline
\texttt{audit} & Audit log writer (PostgreSQL + file fallback) \\
\hline
\texttt{notifier} & Gửi cảnh báo qua Slack, SMTP, Webhook (async) \\
\hline
\texttt{models} & Repository pattern cho \texttt{users}, \texttt{audit\_logs}, \texttt{ip\_lists} \\
\hline
\texttt{database} & Kết nối PostgreSQL, connection pool, migration runner \\
\hline
\end{tabular}
\end{table}

Các package giao tiếp qua interface, không qua struct cụ thể. Đặc biệt, các
interface quan trọng là \texttt{MLPredictor} (trong \texttt{engine}), được
implement bởi \texttt{ml.Client} bên ngoài -- thiết kế này tách hoàn toàn
engine khỏi runtime ML, cho phép unit test engine với mock predictor mà
không cần ML service chạy thật. Chi tiết của thiết kế này được trình bày
ở Chương \ref{section:5.4}.

% ====================================================================
\section{Thiết kế và triển khai Rule Engine}
\label{section:4.2}

\subsection{Cấu trúc dữ liệu Rule}
\label{subsection:4.2.1}

Mỗi rule trong runtime được biểu diễn bằng struct \texttt{Rule} (mô tả ở
file \texttt{internal/engine/types.go}) với các trường chính:

\begin{itemize}[leftmargin=*]
    \item \textbf{Metadata}: \texttt{ID} (định danh duy nhất, ví dụ
    \texttt{SQLI\_001}), \texttt{Category} (\texttt{sqli}, \texttt{xss},
    v.v.), \texttt{Severity} (CRITICAL/HIGH/MEDIUM/LOW), \texttt{Description},
    \texttt{Enabled} (cho phép bật/tắt runtime).

    \item \textbf{Điều kiện áp dụng (\texttt{Conditions})}: chỉ định các
    trường của request mà rule áp dụng -- \texttt{path}, \texttt{query.\{name\}},
    \texttt{headers.\{name\}}, \texttt{body}, \texttt{method}.

    \item \textbf{Transform chain}: danh sách các biến đổi áp dụng tuần tự
    trên text trước khi so khớp pattern.

    \item \textbf{Patterns}: mảng pattern, mỗi pattern có \texttt{Type}
    (\texttt{regex} hoặc \texttt{token}) và payload tương ứng. Với regex,
    pattern là biểu thức RE2; với token, là dãy chuỗi cùng tham số
    \texttt{MaxProximity} (khoảng cách tối đa giữa các token liên tiếp).

    \item \textbf{Scoring}: \texttt{AnomalyScore} (điểm cộng khi rule khớp),
    \texttt{SeverityMultiplier} (hệ số nhân theo severity).

    \item \textbf{Actions}: \texttt{LogOnly} (chỉ ghi log, không cộng điểm),
    \texttt{DenyImmediate} (chặn ngay không qua scoring -- chỉ cho rule
    critical), \texttt{Score} (mặc định, cộng vào tổng anomaly).

    \item \textbf{Exceptions}: danh sách path/header/IP bỏ qua rule này
    (ví dụ \texttt{/health}, \texttt{/metrics} không áp dụng các rule chậm).
\end{itemize}

Khi nạp tập rule từ file JSON \texttt{configs/rules/all\_rules.json}, mỗi
regex được compile một lần và lưu trong struct -- compile-time check đảm
bảo các pattern không hợp lệ bị phát hiện ngay tại load. Tập rule sau khi
parse và compile thành công được hoán đổi (atomic swap) với tập đang chạy,
cho phép \textit{hot reload} không gián đoạn dịch vụ.

Biểu đồ lớp của package \texttt{internal/engine} ở Hình \ref{fig:class-rule-engine}
mô tả quan hệ giữa các kiểu dữ liệu cốt lõi: \texttt{Rule},
\texttt{Selector}, \texttt{Transform}, \texttt{Detect}, \texttt{Engine},
\texttt{Loader} và đặc biệt là interface \texttt{MLPredictor} được khai
báo bên phía consumer để phá bỏ chu trình import sang package
\texttt{internal/ml}.

\begin{figure}[h!]
\centering
\includegraphics[width=\textwidth]{Hinhve/class_rule_engine.png}
\caption{Biểu đồ lớp -- package \texttt{internal/engine} và interface \texttt{MLPredictor}}
\label{fig:class-rule-engine}
\end{figure}

\subsection{Transform chain}
\label{subsection:4.2.2}

Transform chain được hiện thực hóa trong \texttt{internal/engine/transforms.go}
như một map tên-hàm. Mỗi transform là một hàm \texttt{func(string) string}
thuần (no side effect), cho phép áp dụng theo thứ tự một cách dễ dự đoán.
Danh sách transform được hỗ trợ và mô tả ngắn:

\begin{itemize}[leftmargin=*]
    \item \texttt{URL\_DECODE}: giải mã URL một lớp. Lý do chỉ một lớp: tránh
    diễn giải sai dữ liệu hợp pháp chứa ký tự \texttt{\%XX} có ý nghĩa
    (ví dụ password). Đối với double-encoding, rule có thể chèn
    \texttt{URL\_DECODE} hai lần trong chain.

    \item \texttt{LOWERCASE}: chuyển toàn bộ về chữ thường. Cài đặt sử dụng
    \texttt{strings.ToLower} của thư viện chuẩn, hỗ trợ Unicode.

    \item \texttt{COMPRESS\_WHITESPACE}: thay mọi chuỗi \texttt{\textbackslash s+}
    bằng một dấu cách đơn -- chống chèn whitespace giữa các token đặc trưng.

    \item \texttt{HTML\_ENTITY\_DECODE}: giải mã entity HTML
    (\texttt{\&amp;}, \texttt{\&\#x3c;}). Hiện thực qua
    \texttt{html.UnescapeString}.

    \item \texttt{BASE64\_DECODE}: dùng cho rule Log4Shell, nơi JNDI URL
    đôi khi được encode base64. Lưu ý: nếu input không phải base64 hợp lệ,
    transform trả về input gốc thay vì lỗi.

    \item \texttt{REMOVE\_COMMENTS}: xóa các comment SQL (\texttt{/* */},
    \texttt{--}) và shell (\texttt{\#}) -- chống evasion qua chèn comment
    giữa từ khóa.
\end{itemize}

Quan sát thực nghiệm cho thấy chain mặc định \texttt{URL\_DECODE \(\to\)
LOWERCASE \(\to\) COMPRESS\_WHITESPACE} cover khoảng 80\% các kỹ thuật
evasion phổ biến. Các transform còn lại được dùng có chọn lọc cho từng rule
theo đặc thù tấn công.

\subsection{Pattern matching}
\label{subsection:4.2.3}

\textbf{REGEX matching} sử dụng \texttt{regexp.Compile} của Go (engine
RE2). Đặc tính của RE2 là không hỗ trợ backreference và lookaround, đổi
lại đảm bảo độ phức tạp tuyến tính theo độ dài input. Đây là lựa chọn
an toàn đối với WAF: attacker không thể gửi input gây catastrophic
backtracking để DDoS engine -- một lớp tấn công gọi là ReDoS.

\textbf{TOKEN matching} là cơ chế đề tài tự thiết kế, hỗ trợ proximity
match. Cho một dãy token \(t_1, t_2, \ldots, t_n\) và tham số
\texttt{MaxProximity}, hàm match tìm vị trí xuất hiện của \(t_1\), rồi
tìm \(t_2\) trong khoảng \texttt{MaxProximity} ký tự sau đó, và tiếp tục.
Nếu mọi token được tìm thấy theo đúng thứ tự, match thành công. Token match
mạnh hơn regex trong việc bắt các tấn công kiểu \texttt{union(\textbackslash s|\textbackslash /\textbackslash *)*select}
mà không phải viết regex phức tạp dễ sai.

Cả hai kiểu pattern đều hỗ trợ flag \texttt{Negate} (rule khớp khi pattern
không match -- hữu ích cho whitelist-style rule).

\subsection{Anomaly scoring và quyết định}
\label{subsection:4.2.4}

Quy trình tính điểm và quyết định được trình bày bằng pseudocode dưới đây:

\begin{figure}[h!]
\centering
\begin{lstlisting}[style=thesis,numbers=none,basicstyle=\ttfamily\footnotesize]
function Evaluate(request):
    parsed   = Parse(request)
    score    = 0.0
    matched  = []

    for rule in enabled_rules:
        if not applicable(rule, parsed): continue
        text = ApplyTransforms(rule.transforms, field(parsed, rule.cond))
        if MatchPatterns(rule.patterns, text):
            if rule.action == DENY_IMMEDIATE:
                return Decision(BLOCK, rule.score, [rule])
            score += rule.AnomalyScore * rule.SeverityMultiplier
            matched.append(rule)

    if 3.0 <= score < 5.0:                       # gray zone
        ml_result = MLConfirm(parsed, timeout=800ms)
        if ml_result.IsAttack(conf >= 0.8):   score += 2.0
        elif ml_result.IsNormal(conf >= 0.8): score -= 2.0
        # else: leave score unchanged (fail-open)

    return DecisionByThresholds(score, matched)

function DecisionByThresholds(score, matched):
    if score >=  5.0: return BLOCK
    if score >=  4.0: return CHALLENGE
    if score >=  3.0: return LOG
    return ALLOW
\end{lstlisting}
\caption{Pseudocode cơ chế scoring và quyết định}
\label{fig:scoring-pseudo}
\end{figure}

Ba ngưỡng (log/challenge/block) có thể được điều chỉnh runtime qua trang
System Settings. Việc thay đổi ngưỡng có hiệu lực ngay với các request tiếp
theo, không cần restart.

\subsection{Bộ luật 78 rules}
\label{subsection:4.2.5}

78 rule được trình bày dưới dạng JSON trong file \texttt{configs/rules/all\_rules.json}
(toàn văn được liệt kê ở Phụ lục B). Mỗi rule đã được kiểm thử thủ công
trên ít nhất hai biến thể (canonical và một biến thể evasion) trước khi
được đưa vào tập. Bảng \ref{table:sample-rules} trình bày ba ví dụ tiêu
biểu để minh họa cấu trúc.

\begin{table}[h!]
\centering
\caption{Ba ví dụ tiêu biểu: hai chữ ký mạnh và một dấu hiệu yếu (indicator)}
\label{table:sample-rules}
\small
\begin{tabular}{|p{3.1cm}|l|l|p{4.6cm}|}
\hline
\textbf{ID} & \textbf{Nhóm} & \textbf{Severity, Điểm} & \textbf{Pattern \& Transform} \\
\hline
\texttt{WAF-001-SQLI-UNION} & sqli & critical, \(5\times2.0=10\) &
REGEX: \texttt{union} \ldots{} \texttt{select} (cách nhau \(\le 40\) ký tự).\par
Transform: url\_decode \(\to\) lowercase \(\to\) compress\_ws \(\to\) replace\_comments \\
\hline
\texttt{WAF-011-XSS-EVENT-HANDLER} & xss & high, \(4\times1.5=6\) &
REGEX: \texttt{on(error|load|click|\ldots)=} (inline event handler).\par
Transform: url\_decode \(\to\) lowercase \(\to\) html\_decode \\
\hline
\texttt{IND-XSS-002-SCRIPT-WORD} & xss & low, \(2\times0.5=1\) &
REGEX: \texttt{\textbackslash bscript\textbackslash b} (từ \texttt{script} độc lập) -- dấu hiệu yếu, cần tích lũy.\par
Transform: url\_decode \(\to\) html\_decode \(\to\) lowercase \\
\hline
\end{tabular}
\end{table}

% ====================================================================
Tập 78 rule chia làm hai tầng. Khoảng 45 \textit{chữ ký} (signature) bắt các
mẫu tấn công hoàn chỉnh: mỗi chữ ký critical/high thường tự vượt ngưỡng block
khi khớp một mình (ví dụ \texttt{WAF-001} đạt \(10 \ge 5.0\)). Bên cạnh đó, 33
\textit{dấu hiệu} (indicator, tiền tố \texttt{IND-*}) là các tín hiệu yếu
(\texttt{score}\(=2\); severity low/medium) chỉ cộng \(1.0\)--\(2.0\) điểm mỗi
lần khớp. Một dấu hiệu đơn lẻ không đủ để chặn, nhưng nhiều dấu hiệu tích lũy
(ví dụ ba dấu hiệu medium \(=6.0\)) sẽ vượt ngưỡng; còn các request ở mức điểm
giữa \([3.0, 5.0)\) được chuyển sang lớp ML phân xử -- đúng tinh thần
\textit{anomaly scoring} tích lũy của OWASP CRS.

\section{Thiết kế và triển khai ML Layer}
\label{section:4.3}

\subsection{Kiến trúc ML Inference Service}
\label{subsection:4.3.1}

ML Inference Service là một microservice Python độc lập, được triển khai
trong thư mục \texttt{ml-service/} của repo. Service được khởi tạo bằng
FastAPI và Uvicorn (mục \ref{section:3.4}); model checkpoint được nạp vào
RAM ở thời điểm khởi động và giữ nguyên trong suốt vòng đời tiến trình.

Service cung cấp ba endpoint:
\begin{itemize}[leftmargin=*]
    \item \texttt{GET /health}: trả về trạng thái sẵn sàng và phiên bản
    model đang chạy.
    \item \texttt{POST /predict}: nhận một canonical text, trả về
    \texttt{(label, confidence, latency\_ms)}.
    \item \texttt{POST /predict\_batch}: nhận một mảng canonical text, trả
    về mảng tương ứng. Được sử dụng cho mục đích đánh giá hàng loạt và
    không nằm trên hot path của WAF.
\end{itemize}

Logic suy luận bao gồm: tokenize đầu vào (DistilBERT tokenizer, max length
512), forward pass qua model trong context \texttt{torch.no\_grad()},
softmax trên logit để ra phân phối xác suất, lấy argmax cho nhãn và giá trị
xác suất tương ứng cho confidence.

\subsection{Cơ chế tích hợp vào Rule Engine (\texttt{ml\_hook.go})}
\label{subsection:4.3.2}

Tích hợp giữa engine Go và ML service Python được thực hiện qua một file
duy nhất: \texttt{internal/engine/ml\_hook.go}. File này định nghĩa
interface \texttt{MLPredictor} mà engine sử dụng:

\begin{verbatim}
type MLPredictor interface {
    Enabled() bool
    Predict(ctx context.Context, text string) (MLResult, error)
}
\end{verbatim}

Bên trong engine, hàm \texttt{runMLConfirm} kiểm tra \texttt{Enabled()}
trước, sau đó gọi \texttt{Predict} với context có deadline 800 mili giây.
Cấu trúc \texttt{MLResult} chứa \texttt{Label} (string), \texttt{Confidence}
(\(\in [0, 1]\)), và \texttt{LatencyMs} (đo để metric).

Lớp adapter ở phía \texttt{internal/ml/} implement interface trên bằng
cách gửi HTTP POST đến \texttt{http://ml-service:8000/predict} với payload
JSON \texttt{\{"text": canonical\}}. Nếu HTTP client trả lỗi (timeout,
connection refused, mã trả về phi-2xx), \texttt{Predict} trả về
\texttt{MLResult\{\}} kèm error -- engine khi nhận error sẽ giữ nguyên
điểm rule (fail-open).

\subsection{Canonical text preprocessing}
\label{subsection:4.3.3}

Một trong những thách thức kỹ thuật phức tạp nhất của lớp ML là đảm bảo
\textit{Go và Python sinh ra cùng một chuỗi byte cho cùng một request}.
Bất kỳ sự khác biệt nào -- thiếu một dấu xuống dòng, thứ tự header khác
nhau, encode khác nhau -- đều có thể khiến mô hình hoạt động trên một phân
phối khác với phân phối được fine-tune.

Để đảm bảo điều này, đề tài định nghĩa hàm \texttt{BuildCanonicalText}
(Go) và \texttt{canonical\_compose} (Python) với các quy tắc:

\begin{enumerate}[leftmargin=*]
    \item Dòng đầu: \texttt{METHOD PATH}, trong đó PATH đã được URL-decode
    một lớp; nếu có query, append \texttt{?query} với query cũng đã decode.
    \item Các dòng header tiếp theo, mỗi dòng \texttt{Name: Value}, với
    Name viết hoa chữ đầu mỗi từ (canonical-case như HTTP/1.1 spec). Header
    được sắp xếp alphabetic theo Name.
    \item Mỗi value bị cắt còn 256 ký tự nếu dài hơn -- ngăn payload nhồi
    nhét vô tận và đảm bảo input vừa với 512 token của DistilBERT.
    \item Một dòng trống ngăn cách giữa header và body.
    \item Body (nếu có) được append nguyên trạng sau khi cắt còn tối đa
    8192 byte.
\end{enumerate}

Để bảo đảm hai cài đặt khớp nhau, đề tài duy trì một bộ \textbf{7 fixture
test case} chia sẻ giữa hai phía. Mỗi fixture là một cặp \texttt{(raw\_request,
expected\_canonical)}. Pipeline CI kiểm tra cả hai hàm Go và Python trên
cùng các fixture; bất kỳ thay đổi nào ở một phía mà không đồng bộ với phía
kia sẽ làm fail build.

\subsection{Pipeline huấn luyện và phiên bản model}
\label{subsection:4.3.4}

\textbf{Tập dữ liệu huấn luyện}: Được tổng hợp từ ba nguồn. Nguồn thứ nhất
là tập tấn công OWASP CRS payload list và PayloadAllTheThings (khoảng 3000
mẫu tấn công thực). Nguồn thứ hai là synthetic data sinh từ template (ví
dụ template \texttt{SELECT * FROM \{table\} WHERE \{column\} = '\{payload\}'}
với payload là SQLi pattern -- tạo ra hàng nghìn biến thể với độ ngẫu
nhiên kiểm soát). Nguồn thứ ba là traffic ``normal'' được capture từ Juice
Shop và DVWA trong quá trình sử dụng bình thường, cùng với traffic được
sinh bởi tool \texttt{ab} (Apache Bench).

\textbf{Tiền xử lý}: Mọi mẫu đều đi qua hàm \texttt{canonical\_compose}
trước khi cấp cho tokenizer. Điều này đảm bảo phân phối huấn luyện giống
hệt với phân phối inference. Sau tokenize, padding/truncate về độ dài cố
định 512.

\textbf{Quá trình huấn luyện}: Optimizer AdamW, learning rate
\(2 \times 10^{-5}\), weight decay 0.01, batch size 32, max epoch 6 với
early stopping theo validation F1. Loss là cross-entropy với class weights
tỷ lệ nghịch tần suất lớp -- giúp các lớp hiếm như \texttt{log4shell} và
\texttt{ssti} không bị model bỏ qua.

\textbf{Phiên bản model}: Đề tài duy trì ba phiên bản và lưu kết quả huấn
luyện chi tiết cho từng phiên bản:

\begin{itemize}[leftmargin=*]
    \item \textbf{v5 (baseline)}: huấn luyện đầu tiên, đạt accuracy 0.91 trên
    test set nội bộ. \textbf{Đang dùng trong production} do là phiên bản
    cao nhất hiện đạt gate yêu cầu.

    \item \textbf{v6 (fail report)}: cải tiến dataset với một số augmentation
    chung. Accuracy giảm còn 0.8866 do augmentation chưa cân nhắc kỹ phân
    bố lớp. \textbf{Không pass hard gate} (yêu cầu \(\ge 0.92\)), không
    được deploy.

    \item \textbf{v6.1 (kế hoạch)}: 1800 mẫu augment có mục tiêu, tập trung
    vào ba failure pattern đã được phân tích từ v6 -- JSON-body bias, shell
    URL paths, log4shell obfuscation. Hard gate vẫn áp dụng nguyên.
\end{itemize}

Chi tiết của hard gate và phân tích failure được trình bày ở Chương
\ref{section:5.3}.

% ====================================================================
\section{Thiết kế cơ sở dữ liệu}
\label{section:4.4}

\subsection{Sơ đồ ERD}
\label{subsection:4.4.1}

Cơ sở dữ liệu chứa năm bảng chính, được tổ chức theo nguyên tắc tách biệt
quan tâm: \texttt{users} (định danh người dùng nội bộ), \texttt{audit\_logs}
(nhật ký kiểm toán), \texttt{ip\_lists} (danh sách IP), \texttt{rules\_state}
(metadata runtime của rule, không phải định nghĩa rule), và
\texttt{runtime\_config} (cấu hình động không có trong \texttt{config.yaml}).
File migration được liệt kê trong thư mục \texttt{migrations/} và áp dụng
tuần tự khi khởi động.

\begin{figure}[h!]
\centering
\includegraphics[width=\textwidth]{Hinhve/erd.png}
\caption{Sơ đồ ERD của cơ sở dữ liệu WAF}
\label{fig:erd}
\end{figure}

\subsection{Bảng \texttt{audit\_logs}}
\label{subsection:4.4.2}

Bảng \texttt{audit\_logs} có hai loại bản ghi phân biệt qua cột
\texttt{event\_type}: \texttt{request} cho mỗi request được flag (score
\(\ge\) log threshold) và \texttt{system} cho các thao tác quản trị nhạy
cảm. Các cột chi tiết:

\begin{itemize}[leftmargin=*]
    \item Định danh và thời gian: \texttt{id BIGSERIAL PRIMARY KEY},
    \texttt{timestamp TIMESTAMPTZ DEFAULT now()}.
    \item Mạng: \texttt{ip INET}, \texttt{method VARCHAR(10)},
    \texttt{path TEXT}.
    \item Quyết định: \texttt{score NUMERIC(6,2)},
    \texttt{decision VARCHAR(16)} (ENUM-like với CHECK).
    \item Rule match: \texttt{rule\_ids TEXT[]} -- mảng các rule ID đã khớp.
    \item ML output: \texttt{ml\_label VARCHAR(32)}, \texttt{ml\_confidence
    NUMERIC(4,3)}.
    \item Liên kết người dùng: \texttt{actor\_id INTEGER REFERENCES
    users(id)} -- NULL với event\_type = request, không NULL với system.
    \item Phân loại: \texttt{event\_type VARCHAR(16)} với CHECK
    \texttt{IN ('request','system')}.
\end{itemize}

Chỉ mục B-tree trên \texttt{timestamp} và \texttt{ip} đảm bảo các truy
vấn lọc theo thời gian hoặc IP phản hồi dưới 100 mili giây ngay cả với
bảng nhiều triệu bản ghi.

\subsection{Bảng \texttt{users}}
\label{subsection:4.4.3}

Bảng \texttt{users} là trái tim của lớp xác thực và phân quyền. Cấu trúc
hoàn chỉnh:

\begin{verbatim}
CREATE TABLE users (
    id            SERIAL       PRIMARY KEY,
    username      VARCHAR(64)  UNIQUE NOT NULL,
    email         VARCHAR(255) UNIQUE,
    password_hash VARCHAR(255) NOT NULL,    -- bcrypt cost 10
    role          VARCHAR(16)  NOT NULL
                  CHECK (role IN ('admin','editor','viewer')),
    created_at    TIMESTAMPTZ DEFAULT now(),
    updated_at    TIMESTAMPTZ DEFAULT now()
);

CREATE TRIGGER update_users_updated_at
BEFORE UPDATE ON users
FOR EACH ROW EXECUTE FUNCTION set_updated_at();
\end{verbatim}

Ba điểm thiết kế quan trọng. Thứ nhất, ràng buộc \texttt{CHECK} trên
\texttt{role} chuyển trách nhiệm validate role từ tầng ứng dụng xuống tầng
schema -- ngay cả khi có lỗi logic ở handler, không thể insert một role
không hợp lệ vào DB. Thứ hai, \texttt{UNIQUE} trên cả \texttt{username}
và \texttt{email} ngăn trùng tài khoản ở mức nguyên tử (race condition giữa
hai admin tạo đồng thời cùng một username sẽ chỉ một insert thành công).
Thứ ba, trigger tự động cập nhật \texttt{updated\_at} giải phóng tầng
ứng dụng khỏi việc phải nhớ set thủ công.

Quan hệ giữa bảng \texttt{users} và luồng xác thực JWT được biểu diễn
trên biểu đồ tuần tự ở Hình \ref{fig:seq-auth-login}. Hình này phác họa
hai pha: pha đăng nhập (so khớp bcrypt, sinh JWT, đặt cookie
\texttt{HttpOnly}) và pha truy cập tài nguyên kế tiếp (verify JWT, kiểm
tra role theo middleware tương ứng).

\begin{figure}[h!]
\centering
\includegraphics[width=0.9\textwidth,keepaspectratio]{Hinhve/seq_auth_login.png}
\caption{Biểu đồ tuần tự -- Đăng nhập và truy cập tài nguyên qua JWT}
\label{fig:seq-auth-login}
\end{figure}

% ====================================================================
\section{Giao diện Dashboard}
\label{section:4.5}

\subsection{Tổng quan giao diện}
\label{subsection:4.5.1}

Dashboard là một single-page application thuần (HTML/CSS/JavaScript) không
sử dụng framework SPA (React/Vue/Angular). Lựa chọn này có ba lý do:
\textbf{(i)} giảm thiểu bề mặt tấn công của chính dashboard -- không có
build pipeline phức tạp, không có hàng trăm gói \texttt{npm} với phụ thuộc
chuyển tiếp; \textbf{(ii)} bundle assets được nhúng trực tiếp vào binary
WAF qua \texttt{//go:embed} của Go, không cần CDN hay file server riêng;
\textbf{(iii)} giảm độ phức tạp cho người maintain -- mã frontend gần với
spec HTML/CSS gốc, không yêu cầu kiến thức về cú pháp framework.

Để bù lại thiếu hụt về thư viện UI, dashboard sử dụng \textbf{Tailwind CSS
qua CDN} cho styling, kết hợp với một bộ CSS custom properties định nghĩa
theme. Tất cả các trang dùng font \textbf{Inter} (lấy từ Google Fonts) cho
giao diện thống nhất.

\subsection{Các màn hình chính}
\label{subsection:4.5.2}

Dashboard gồm chín trang HTML, được liệt kê trong bảng \ref{table:web-pages}.
Mọi trang đều được phục vụ qua \texttt{//go:embed} từ thư mục \texttt{web/},
không yêu cầu file server bên ngoài.

\begin{table}[h!]
\centering
\caption{Danh sách các trang giao diện}
\label{table:web-pages}
\small
\begin{tabular}{|l|l|p{6cm}|}
\hline
\textbf{Trang} & \textbf{Auth} & \textbf{Chức năng} \\
\hline
\texttt{login.html} & public & Form đăng nhập, JWT set vào HttpOnly cookie \\
\hline
\texttt{register.html} & public & Đăng ký (role bị fix = \texttt{viewer}) \\
\hline
\texttt{index.html} & any & Dashboard chính: counter real-time, top attacks, logs gần nhất, các tab quản lý rule/IP \\
\hline
\texttt{settings.html} & any & Account settings: đổi email, đổi mật khẩu (yêu cầu old password) \\
\hline
\texttt{users.html} & admin & User management: tạo/sửa/xóa user, reset password, ràng buộc invariant \\
\hline
\texttt{waf\_config\_ui.html} & admin/editor & System Settings: ngưỡng quyết định, URL backend, ML endpoint \\
\hline
\texttt{rule\_builder.html} & admin/editor & Editor JSON nâng cao để tạo rule mới \\
\hline
\end{tabular}
\end{table}

Mỗi trang gọi WAF API qua \texttt{fetch()} với cookie kèm theo tự động.
Đối với các thao tác state-changing (POST/PUT/DELETE), client kiểm tra
response status và hiển thị thông báo lỗi tương ứng. Các trường nhập liệu
được validate ở cả client (UX) và server (security) -- nguyên tắc
\textit{double validation}.

\subsection{Hệ thống Theme đồng nhất}
\label{subsection:4.5.3}

Một trong những hạng mục cải tiến cuối được thực hiện là thống nhất theme
giữa tất cả các trang. Trước đó, \texttt{login.html} dùng gradient tím-xanh
hard-coded, trong khi dashboard chính dùng theme đen tối hiện đại -- gây
cảm giác hai sản phẩm khác nhau khi user di chuyển giữa hai trang.

Giải pháp gồm bốn thành phần:

\begin{enumerate}[leftmargin=*]
    \item \textbf{CSS custom properties với \texttt{data-theme} attribute}:
    Mọi màu sắc, background, border được khai báo dưới dạng biến CSS
    (\texttt{--bg-primary}, \texttt{--text-primary}, v.v.) và được override
    cho \texttt{[data-theme="light"]}. Việc đổi theme chỉ cần thay attribute
    \texttt{data-theme} trên \texttt{<html>}.

    \item \textbf{Pre-paint bootstrap script chống FOUC}: Một đoạn script
    inline ở \texttt{<head>} đọc \texttt{localStorage.waf\_theme} và set
    attribute ngay trước khi render -- tránh ``flash of unstyled content''
    khi user reload trang.

    \item \textbf{Nút toggle theme}: Mỗi trang có nút ở góc trên-phải cho
    phép chuyển dark/light. Lựa chọn được lưu vào \texttt{localStorage}
    để các trang khác cũng đồng bộ.

    \item \textbf{Glass-card layout thống nhất}: Mọi form và panel được
    bọc trong CSS class \texttt{.glass-card} với backdrop-filter và border
    cùng tone -- đồng nhất visual identity qua các trang.
\end{enumerate}

% ====================================================================
\section{Kiểm thử và Đánh giá}
\label{section:4.6}

\subsection{Chiến lược kiểm thử}
\label{subsection:4.6.1}

Đề tài áp dụng chiến lược kiểm thử bốn lớp:

\begin{itemize}[leftmargin=*]
    \item \textbf{Unit test}: Kiểm thử các hàm thuần (transforms, pattern
    matching, scoring) bằng \texttt{go test}. Coverage mục tiêu \(\ge\) 70\%
    cho các package core (\texttt{engine}, \texttt{auth}, \texttt{models}).

    \item \textbf{Integration test}: Khởi động WAF + PostgreSQL + ML
    service trong Docker Compose, kiểm thử các kịch bản end-to-end (đăng
    nhập, tạo user, gửi request tấn công, xem log).

    \item \textbf{Penetration test (security test)}: Phát 42 payload tấn
    công OWASP chuẩn (SQLi, XSS, CMDi, Path Traversal, SSRF, XXE, Log4Shell,
    SSTI, NoSQLi, RCE) qua WAF, đo block rate và phân loại các trường
    hợp bypass.

    \item \textbf{Load test}: Apache Bench gửi 1000 request, concurrency
    100, đo throughput và latency P50/P95/P99.
\end{itemize}

\subsection{Kết quả kiểm thử Rule Engine}
\label{subsection:4.6.2}

Lớp rule được đánh giá bằng harness tái lập \texttt{cmd/wafeval} (mô tả ở mục
\ref{subsection:5.2.3}): nạp trực tiếp 78 rule production qua package
\texttt{internal/engine} với đúng ngưỡng \texttt{config.yaml}, chạy 46 payload
tấn công (10 nhóm OWASP) và 52 request hợp lệ cố tình gây nhiễu, lấy điểm rule
làm tín hiệu quyết định (không gọi ML). Bảng \ref{table:rule-test-results} tổng
kết.

\begin{table}[h!]
\centering
\caption{Kết quả kiểm thử rule engine (78 rule, không kèm ML)}
\label{table:rule-test-results}
\small
\begin{tabular}{|l|c|}
\hline
\textbf{Tiêu chí} & \textbf{Kết quả} \\
\hline
Rules nạp thành công vào engine & \textbf{78/78} \\
\hline
Block rate (46 payload OWASP) & \textbf{87.0\% (40/46)} \\
\hline
Caught rate (block + challenge) & \textbf{89.1\% (41/46)} \\
\hline
Nhóm injection chính (SQLi/XSS/LFI/XXE/NoSQLi) & \textbf{100\% block} \\
\hline
False positive (52 request đối nghịch) & \textbf{3.8\% (2/52)} \\
\hline
Latency rule engine (mỗi request) & \textbf{sub-ms (\(\approx 0.13\) ms)} \\
\hline
\end{tabular}
\end{table}

Năm payload lọt khỏi lớp rule đều thuộc các trường hợp \textit{cần lớp ML hoặc
cơ chế stateful}: SSTI cô lập \texttt{\{\{7*7\}\}} (chỉ một dấu hiệu indicator),
SSRF IPv6 loopback \texttt{http://[::1]/}, UA credential-stuffing, probe file
backup, và ATO một-lần (chỉ kích hoạt khi đếm nhiều lần thử). Đây đúng là vùng
được thiết kế để dấu hiệu tích lũy và lớp ML xác minh đảm nhận (phân tích chi
tiết ở mục \ref{subsection:5.2.3}).

\subsection{Kết quả kiểm thử ML Model}
\label{subsection:4.6.3}

Mô hình v5 (đang dùng trong production) được đánh giá trên tập 97 mẫu
out-of-distribution (OOD) -- không nằm trong tập huấn luyện. Kết quả tổng
hợp ở mức binary attack/normal:

\begin{itemize}[leftmargin=*]
    \item Precision (attack): 0.976
    \item Recall (attack): 0.988
    \item F1 (attack): 0.982
    \item Accuracy chung: 0.928
\end{itemize}

Phân tích 11 trường hợp misclassification cho thấy ba pattern lỗi tập
trung: \textbf{(i)} \textit{JSON-body bias} -- model có xu hướng dự đoán
\texttt{normal} cho các request có body là JSON dù chứa payload SQLi
(do mất cân bằng dữ liệu: phần lớn JSON body trong tập train là benign);
\textbf{(ii)} \textit{Shell URL path} -- một số RCE payload trong path
không được nhận diện do path quá ngắn so với chiều dài trung bình mà
model học; \textbf{(iii)} \textit{Log4Shell obfuscation} -- các biến
thể \texttt{\$\{\$\{lower:j\}ndi:} không trùng pattern trong tập train.
Ba pattern lỗi này định hướng cho việc augmentation ở v6.1.

\subsection{Kết quả kiểm thử Rate Limiting}
\label{subsection:4.6.4}

Kịch bản kiểm thử: gửi burst 1500 request từ cùng một IP với rate giả
lập, cấu hình rate limiter \texttt{rate = 10 req/s, burst = 50}. Kết quả:

\begin{itemize}[leftmargin=*]
    \item Block rate: 984/1500 (\textbf{65.6\%}). Trong 516 request lọt
    qua, 50 request đầu vào trong burst window, 466 request còn lại thuộc
    về các time slice tiếp theo với rate được phân bố lại.
    \item Latency trung bình thêm vào do rate limiter: \(< 50\,\mu\)s.
    \item Không có rò rỉ bộ nhớ sau 24 giờ chạy với 100K IP unique.
\end{itemize}

Số liệu này được coi là chấp nhận được với mục tiêu cản trở brute-force và
DDoS ứng dụng quy mô vừa. Để chống DDoS volumetric, cần một lớp phía
trước (NGINX rate limit, CDN, hoặc dịch vụ chuyên dụng).

\subsection{Kiểm thử bảo mật nội bộ (auth/authorization)}
\label{subsection:4.6.5}

Đây là phần kiểm thử quan trọng nhất đối với lớp xác thực, vì các lỗ hổng
ở đây có khả năng vô hiệu hóa toàn bộ phần còn lại của hệ thống. Bảng
\ref{table:auth-tests} tổng kết các kịch bản đã kiểm thử.

\begin{table}[h!]
\centering
\caption{Kết quả kiểm thử bảo mật lớp auth}
\label{table:auth-tests}
\small
\begin{tabular}{|p{6.5cm}|l|c|}
\hline
\textbf{Kịch bản test} & \textbf{Mã kỳ vọng} & \textbf{Pass?} \\
\hline
Gọi admin endpoint không có JWT & 401 & \checkmark \\
\hline
Gọi admin endpoint với viewer JWT & 403 & \checkmark \\
\hline
Gọi admin endpoint với admin JWT & 200 & \checkmark \\
\hline
\texttt{POST /register \{"role":"admin"\}} & 200, role=viewer & \checkmark \\
\hline
Xóa admin cuối cùng & 400 \texttt{"last admin"} & \checkmark \\
\hline
Demote admin cuối cùng từ admin sang viewer & 400 \texttt{"last admin"} & \checkmark \\
\hline
Admin tự xóa chính mình & 400 \texttt{"own account"} & \checkmark \\
\hline
Admin tự demote chính mình & 400 \texttt{"own account"} & \checkmark \\
\hline
Đổi password với old\_password sai & 401 & \checkmark \\
\hline
Đổi password với old\_password đúng & 200 + audit log & \checkmark \\
\hline
\end{tabular}
\end{table}

Đặc biệt quan trọng là kịch bản test regression \texttt{POST /register
\{"role":"admin"\}}: trước fix, request này tạo được tài khoản admin từ
public endpoint -- lỗ leo thang đặc quyền nghiêm trọng. Sau fix, trường
\texttt{role} bị bỏ qua hoàn toàn ở handler và role luôn được hard-code
\texttt{viewer}. Việc tạo tài khoản với role cao hơn chỉ được thực hiện
qua endpoint admin \texttt{POST /waf-api/auth/users}.

% ====================================================================
\section{Triển khai}
\label{section:4.7}

\subsection{Yêu cầu hệ thống}
\label{subsection:4.7.1}

Yêu cầu phần cứng tối thiểu để chạy đầy đủ stack (WAF + PostgreSQL + ML
service):

\begin{itemize}[leftmargin=*]
    \item CPU: 2 vCPU (x86\_64 hoặc ARM64).
    \item RAM: 2 GB cho WAF core và PostgreSQL; thêm 1.5 GB cho ML service
    với model v5 (model + tokenizer + PyTorch runtime).
    \item Đĩa: 5 GB cho image Docker và 10 GB cho dữ liệu PostgreSQL (audit
    log với 90 ngày lưu giữ ở mức traffic 100 req/s).
    \item Phần mềm: Docker \(\ge\) 20.10, Docker Compose \(\ge\) 2.0,
    Linux kernel \(\ge\) 4.18 (cho cgroup v2).
\end{itemize}

Trên các môi trường không đủ RAM, ML service có thể được tắt -- WAF sẽ
hoạt động chỉ với rule engine. Cấu hình này được kiểm chứng vẫn đạt block
rate \(\ge 95\)\% trên 42 OWASP payload chuẩn (mất khả năng xử lý vùng
xám nhưng không ảnh hưởng các trường hợp rõ ràng).

Mô hình triển khai chuẩn của hệ thống được mô tả trong Hình
\ref{fig:deployment}. Các thành phần WAF, ML service, PostgreSQL và
upstream application có thể đặt trên cùng một host (mô hình single-node
cho môi trường thử nghiệm) hoặc tách thành nhiều host vật lý -- tiến
trình WAF chỉ giao tiếp với các phụ trợ qua mạng nội bộ (loopback hoặc
VLAN private), không qua Internet.

\begin{figure}[h!]
\centering
\includegraphics[width=\textwidth]{Hinhve/deployment.png}
\caption{Biểu đồ triển khai -- Phân bổ các thành phần lên host}
\label{fig:deployment}
\end{figure}

\subsection{Hướng dẫn triển khai (Docker Compose)}
\label{subsection:4.7.2}

Quy trình triển khai gồm năm bước, mỗi bước là một lệnh đơn:

\begin{verbatim}
# Bước 1: clone repo và cd vào thư mục
git clone <repo-url> && cd waf-project

# Bước 2: khởi tạo PostgreSQL và chạy migration
./scripts/setup_db.sh

# Bước 3: sinh self-signed TLS certificate cho cổng 8443
./scripts/generate_certs.sh

# Bước 4: build và chạy WAF core (cổng 8443)
make build && make run

# Bước 5: khởi động ML service (cổng 8000)
docker compose up ml -d
\end{verbatim}

Sau khi cả năm bước hoàn tất, mở trình duyệt và truy cập
\texttt{https://localhost:8443/login.html} với tài khoản mặc định
\texttt{admin / admin123}. \textbf{Lưu ý bảo mật quan trọng}: tài khoản
mặc định này phải được đổi mật khẩu ngay sau lần đăng nhập đầu (qua trang
Settings) trong môi trường có khả năng tiếp xúc bên ngoài.

\subsection{Cấu trúc file triển khai}
\label{subsection:4.7.3}

Các file liên quan đến triển khai:

\begin{itemize}[leftmargin=*]
    \item \texttt{docker-compose.db.yml}: định nghĩa hai service PostgreSQL
    và ML, kèm network và volume.

    \item \texttt{Makefile}: cung cấp các target \texttt{build}, \texttt{run},
    \texttt{test}, \texttt{clean}, \texttt{docker-build}, \texttt{docker-push}
    -- abstract hóa các lệnh \texttt{go build} và \texttt{docker} với
    tham số mặc định.

    \item \texttt{configs/config.yaml}: file cấu hình chính, gồm
    \texttt{server}, \texttt{upstream}, \texttt{parser}, \texttt{rate\_limit},
    \texttt{behavior}, \texttt{decision}, \texttt{database}, \texttt{auth},
    \texttt{admin}, \texttt{ml}, \texttt{notifier}.

    \item \texttt{configs/rules/all\_rules.json}: tập 78 rule. Được nạp ở
    thời điểm khởi động và có thể hot-reload qua API admin.

    \item \texttt{scripts/setup\_db.sh}, \texttt{scripts/generate\_certs.sh}:
    script hỗ trợ khởi tạo môi trường.

    \item \texttt{migrations/}: bốn migration SQL (\texttt{001\_create\_users.sql},
    \texttt{002\_fix\_admin\_password.sql}, \texttt{003\_create\_runtime\_config.sql},
    \texttt{004\_create\_runtime\_state.sql}). Migration runner chạy tuần
    tự khi WAF khởi động và idempotent (chạy lại không gây lỗi).
\end{itemize}

\bigskip
\noindent\textbf{Kết chương.}
Chương này đã trình bày toàn bộ quá trình thiết kế và hiện thực hóa hệ
thống WAF, từ kiến trúc tổng thể với chín bước xử lý request, cho đến các
chi tiết về rule engine, ML layer, cơ sở dữ liệu, dashboard và triển khai
thực tế. Kết quả kiểm thử cho thấy lớp rule đạt block rate 87.0\% (bắt được
89.1\%) trên 46 OWASP payload đa nhóm, ML binary precision 0.976 / recall 0.988, và mọi
kịch bản kiểm thử bảo mật lớp auth đều pass -- bao gồm các test regression
chống lỗ leo thang đặc quyền và các invariant bảo vệ chống lockout. Các
phát hiện thú vị trong quá trình thiết kế và đánh giá -- đặc biệt là cơ
chế anomaly scoring hai lớp, transform chain chống evasion, pipeline ML
có hard gate, và thiết kế interface tách biệt engine-ML -- sẽ được phân
tích sâu hơn ở Chương \ref{section:5.1} với tư cách là các đóng góp nổi
bật của đề tài.

\end{document}
